\chapter{Inequalities, Convexity and Estimates}\label{ch:inequalities}

\begin{center}\itshape
The number $e$ appears in an inequality about averages of an arbitrary positive sequence.\\
Why there, and why is it the best possible constant?
\end{center}

Take the two positive numbers $1$ and $4$. Their arithmetic mean is $2.5$. Their geometric mean is $2$. The gap is a square:
\[
\left(\frac{1+4}2\right)^2-\bigl(\sqrt{1\cdot4}\bigr)^2=\left(\frac{4-1}2\right)^2.
\]
That identity is the first inequality of the chapter, and it is sharp: the two means agree precisely when the two numbers agree. The same square, rewritten for a graph, for two lists, and for a series, produces Jensen's inequality, the comparison of means, Cauchy--Schwarz, H\"older, the power means, and Minkowski's triangle inequality. At the end of that chain the constant $e$ appears in front of a series, and no smaller constant works for every series.

\section{A square is never negative}\label{sec:squares}

The basic move is to write the difference of the two sides as something you can see is nonnegative.

\begin{reminder}[moving an inequality]
You may add inequalities: $a\le b$ and $c\le d$ give $a+c\le b+d$. You may multiply by a positive number. Multiplying by a negative number reverses the inequality. If $0<a\le b$ and $0<c\le d$, then $ac\le bd$. On nonnegative numbers, squaring preserves an inequality, and so does taking square roots. On positive numbers, taking reciprocals reverses it. Subtraction needs the bounds the other way round: $a\le b$ and $c\le d$ give $a-d\le b-c$. The reliable proof that $A\le B$ exhibits $B-A$ as a square, a sum of squares, or a product of nonnegative numbers.
\end{reminder}

Start with $(a-b)^2\ge0$, for any real $a$ and $b$. Expand:
\[
a^2-2ab+b^2\ge0,\qquad\text{so}\qquad a^2+b^2\ge2ab.
\]
Equality holds exactly when $a=b$, because that is when the square vanishes. Now let $x$ and $y$ be nonnegative, and put $a=\sqrt x$ and $b=\sqrt y$:
\[
x+y\ge2\sqrt{xy},\qquad\text{so}\qquad\sqrt{xy}\le\frac{x+y}2.
\]
Equality holds exactly when $x=y$. For $x=1$ and $y=4$ this is $2\le2.5$. Figure~\ref{fig:frame} is the same identity drawn as areas, with $a=3$ and $b=1$: the outer square has area $16$, the inner square has area $4$, and the four rectangles have area $12=4\cdot3\cdot1$. Since the inner square is nonnegative,
\[
(a+b)^2\ge4ab,\qquad\frac{a+b}2\ge\sqrt{ab}.
\]

\begin{figure}[ht]
\centering
\fig{ch04-frame}
\caption{The identity $(a+b)^2=4ab+(a-b)^2$ for $a=3$ and $b=1$. The outer square has area $16$. The ochre square in the middle has area $(a-b)^2=4$. Each of the four rectangles has area $ab=3$. What remains after the inner square is removed is $4ab$, so the arithmetic mean $2$ is at least the geometric mean $\sqrt3$.}
\label{fig:frame}
\end{figure}

The same comparison has a geometric form that also sees the harmonic mean. Let the diameter of a semicircle have length $a+b$, with $a>0$ and $b>0$, and mark the point that splits the diameter into segments $a$ and $b$. The radius is $(a+b)/2$. The vertical line at the split meets the semicircle at height $h$. To see why $h=\sqrt{ab}$, put the diameter on the interval from $0$ to $a+b$. The centre is at $(a+b)/2$ and the radius is the same number. The point $(a,h)$ lies on the circle, so
\[
\left(a-\frac{a+b}2\right)^2+h^2=\left(\frac{a+b}2\right)^2.
\]
The first term is $(a-b)^2/4$. Subtract it:
\[
h^2=\frac{(a+b)^2-(a-b)^2}4=ab.
\]
The height is the geometric mean. It cannot exceed the radius, because a point of the semicircle lies at most one radius above the diameter. That is the inequality $\sqrt{ab}\le(a+b)/2$ again.

The harmonic mean is the length of a shadow. Put the diameter on the interval from $0$ to $D=a+b$. The foot of the altitude is $F=(a,0)$, the point on the semicircle is $C=(a,\sqrt{ab})$, and the centre is $M=(D/2,0)$. The vector $v$ from $F$ to $C$ is $(0,\sqrt{ab})$. The vector from $M$ to $C$ is $\bigl((a-b)/2,\sqrt{ab}\bigr)$, of length $D/2$. Divide it by that length. The unit vector in the same direction is
\[
u=\left(\frac{a-b}{D},\frac{2\sqrt{ab}}{D}\right).
\]
The projection of $v$ onto $u$ has length $|v\cdot u|$, and
\[
v\cdot u=\sqrt{ab}\cdot\frac{2\sqrt{ab}}{D}=\frac{2ab}{a+b}.
\]
Write $h=\sqrt{ab}$ and $u=(u_1,u_2)$. Then $v\cdot u=hu_2$, and $u_1^2+u_2^2=1$, so $u_2^2\le1$ and $|v\cdot u|\le h$. Therefore
\[
\frac{2ab}{a+b}\le\sqrt{ab}\le\frac{a+b}2.
\]
Figure~\ref{fig:semicircle} draws $a=1$ and $b=4$. The three lengths are $1.6$, $2$ and $2.5$.

\begin{figure}[ht]
\centering
\fig{ch04-semicircle}
\caption{Three means on a semicircle of diameter $a+b=5$, with $a=1$ and $b=4$. The charcoal radius is the arithmetic mean $2.5$. The rust altitude is the geometric mean $2$. The ochre segment is the harmonic mean $1.6$. The altitude is at most the radius, and the ochre segment is a projection of the altitude, so it is at most the altitude.}
\label{fig:semicircle}
\end{figure}

Add the three pairwise inequalities $2ab\le a^2+b^2$, $2bc\le b^2+c^2$ and $2ca\le c^2+a^2$. You get
\[
ab+bc+ca\le a^2+b^2+c^2,
\]
with equality exactly when $a=b=c$.

\begin{examplex}
For $x>0$, the claim $x+1/x\ge2$ is the square $(x-1)^2\ge0$ after dividing by $x$. Expand the square, add $2x$, and divide by the positive number $x$:
\[
(x-1)^2\ge0,\qquad x^2+1\ge2x,\qquad x+\frac1x\ge2.
\]
Equality holds at $x=1$. The same pairing works for three positive numbers. Expand
\[
(a+b+c)\left(\frac1a+\frac1b+\frac1c\right)=3+\left(\frac ab+\frac ba\right)+\left(\frac bc+\frac cb\right)+\left(\frac ca+\frac ac\right).
\]
Each pair $t+1/t$ is at least $2$, so the sum is at least $3+2+2+2=9$. Equality holds exactly when $a=b$, $b=c$ and $c=a$, that is, when $a=b=c$.
\end{examplex}

Two comparisons from earlier chapters control absolute values. Proposition~\ref{prop:real-triangle} says
\[
|a+b|\le|a|+|b|\qquad\text{and}\qquad\bigl||a|-|b|\bigr|\le|a-b|.
\]
They let you split a difficult difference into pieces you can estimate separately.

Proposition~\ref{prop:bernoulli} is the inequality $(1+x)^n\ge1+nx$ for $x\ge-1$ and integers $n\ge0$, proved there by induction. It says that $(1+x)^n$ grows at least as fast as the straight line $1+nx$. Jacob Bernoulli printed the comparison in $1689$, in his tract on infinite series. It is older than that tract: Isaac Barrow's geometrical lectures of $1670$ already contain it. Exercise~\ref{ex:bernoulli} gets two concrete bounds out of the form proved in this book.

\portrait{jbernoulli}{Jacob Bernoulli (1655--1705). Portrait by his brother Nicolaus Bernoulli, 1687. Public domain, via Wikimedia Commons.}

\section{Chords, tangents and averages}\label{sec:convex}

The square compared two numbers. A chord compares as many as you like, once you know which way the graph bends.

A function $f$, defined on an interval $I$, is \emph{convex} when for all $x,y\in I$ and all $\lambda\in[0,1]$,
\[
f\bigl(\lambda x+(1-\lambda)y\bigr)\le\lambda f(x)+(1-\lambda)f(y).
\]
The left side is the height of the graph at a point of the segment from $x$ to $y$. The right side is the height of the chord at that same point. The graph sits under the chord. A function is \emph{concave} when $-f$ is convex, so the inequality reverses and the graph sits above its chords. The graph of $x^2$ holds water. The graph of $\ln$ spills it.

\begin{proposition}[convexity tests]\label{prop:convex}
Let $f:I\to\R$ be twice differentiable on an interval $I$.
\begin{enumerate}[label=(\alph*)]
\item $f$ is convex if and only if $f''\ge0$ on $I$.
\item If $f$ is convex, then $f$ lies above its tangents: $f(x)\ge f(m)+f'(m)(x-m)$ for all $x,m\in I$. If $f''>0$, this inequality is strict for $x\ne m$.
\end{enumerate}
\end{proposition}
\begin{proof}
\emph{Positive second derivative gives the chord.} Suppose $f''\ge0$. Then $f'$ is non-decreasing: if $u<v$, the mean value theorem (Theorem~\ref{thm:mvt}) supplies $c$ between them with $f'(v)-f'(u)=f''(c)(v-u)\ge0$. Take $x<y$ in $I$ and $\lambda\in\,]0,1[$, and set $z=\lambda x+(1-\lambda)y$. Then $z$ lies between $x$ and $y$, and
\[
z-x=(1-\lambda)(y-x),\qquad y-z=\lambda(y-x).
\]
Apply the mean value theorem on $[x,z]$ and on $[z,y]$. There are $\xi_1\in\,]x,z[$ and $\xi_2\in\,]z,y[$ with
\[
\frac{f(z)-f(x)}{z-x}=f'(\xi_1)\le f'(\xi_2)=\frac{f(y)-f(z)}{y-z},
\]
where the middle comparison uses $\xi_1<\xi_2$ and the fact that $f'$ is non-decreasing. Clear the positive denominators. Multiply the inequality by $(z-x)(y-z)$:
\[
(y-z)\bigl(f(z)-f(x)\bigr)\le(z-x)\bigl(f(y)-f(z)\bigr).
\]
Insert $y-z=\lambda(y-x)$ and $z-x=(1-\lambda)(y-x)$, and cancel the positive factor $y-x$:
\[
\lambda\bigl(f(z)-f(x)\bigr)\le(1-\lambda)\bigl(f(y)-f(z)\bigr).
\]
Expand and collect the terms in $f(z)$:
\[
\lambda f(z)+(1-\lambda)f(z)\le\lambda f(x)+(1-\lambda)f(y),
\]
which is $f(z)\le\lambda f(x)+(1-\lambda)f(y)$. The endpoint cases $\lambda=0$ and $\lambda=1$ are equalities. So $f$ is convex.

\emph{A convex graph forces the second derivative.} Suppose $f$ is convex, and take $x<z<y$. The chord inequality at this $z$, solved for the slopes, is the comparison just run backwards:
\[
\frac{f(z)-f(x)}{z-x}\le\frac{f(y)-f(z)}{y-z}.
\]
The slope inequality holds for every $z$ between $x$ and $y$. Send $z$ down to $x$: the left side tends to $f'(x)$, and the right side tends to $(f(y)-f(x))/(y-x)$. Separately, send $z$ up to $y$: the right side tends to $f'(y)$, and the left side tends to the same difference quotient. Therefore
\[
f'(x)\le\frac{f(y)-f(x)}{y-x}\le f'(y).
\]
So $f'$ is non-decreasing. Its difference quotients are nonnegative, and $f''$ is their limit, so $f''\ge0$.

\emph{The tangent.} Fix $m\in I$ and set $g(x)=f(x)-f(m)-f'(m)(x-m)$. Then $g(m)=0$ and $g'=f'-f'(m)$. Part (a) says $f'$ is non-decreasing, so $g'\le0$ on the left of $m$ and $g'\ge0$ on the right. Thus $g$ decreases down to $m$ and increases afterwards, and $g(x)\ge g(m)=0$. If $f''>0$, then $f'$ is strictly increasing, so $g'$ is strictly negative on the left of $m$ and strictly positive on the right. Hence $g(x)>0$ for $x\ne m$.
\end{proof}

\begin{examplex}
The second-derivative test gives a list at once. Convex on the obvious domains: $x^2$, $e^x$, $x^p$ on $]0,\infty[$ for $p\ge1$ and for $p\le0$, $-\ln x$, $x\ln x$ on $]0,\infty[$ (its second derivative is $1/x$), and $1/x$ on $]0,\infty[$. Concave: $\ln x$, $\sqrt x$, $x^p$ for $0<p<1$, and $\sin$ on $[0,\pi]$. Differentiability is not part of the definition. The absolute value $|x|$ is convex, because a convex combination of absolute values is at least the absolute value of the combination, and it has no derivative at $0$. Every later inequality that uses only the chord definition applies to it.
\end{examplex}

\begin{think}
Draw a convex graph and choose three points $x_1,x_2,x_3$ on the axis. The point
\[
\left(\frac{x_1+x_2+x_3}3,\ \frac{f(x_1)+f(x_2)+f(x_3)}3\right)
\]
is the centre of gravity of the triangle whose vertices lie on the graph. Does that centre sit above the graph, or below it?
\end{think}

\begin{figure}[ht]
\centering
\fig{ch04-jensen}
\caption{Jensen's inequality for three points of a convex graph. Each side of the triangle is a chord, so the triangle lies above the graph. The centre of gravity has horizontal coordinate $\bar x=(x_1+x_2+x_3)/3$ and height $(f(x_1)+f(x_2)+f(x_3))/3$. It lies in the triangle, hence above the point $(\bar x,f(\bar x))$ on the graph.}
\label{fig:jensen}
\end{figure}

Figure~\ref{fig:jensen} answers the question in the box: the centre sits above the graph. The general fact is Jensen's inequality. The weights are allowed to be unequal, as long as they are nonnegative and add to $1$.

\begin{theorem}[Jensen]\label{thm:jensen}
Let $f$ be convex on an interval $I$. Let $x_1,\dots,x_n\in I$ and let $\lambda_1,\dots,\lambda_n\ge0$ with $\sum_{i=1}^n\lambda_i=1$. Then
\[
f\left(\sum_{i=1}^n\lambda_ix_i\right)\le\sum_{i=1}^n\lambda_if(x_i).
\]
\end{theorem}
\begin{proof}
The case $n=1$ is the equality $f(x_1)=f(x_1)$. The case $n=2$ is the definition of convexity. Assume the claim is known for $n-1$ weights, where $n\ge3$, and consider $n$ weights.

If $\lambda_n=1$, then every earlier weight is $0$, and both sides equal $f(x_n)$. If $\lambda_n=0$, the claim is the induction hypothesis on the first $n-1$ weights. Now suppose $0<\lambda_n<1$, and set $\mu=1-\lambda_n$. Then $\mu\in\,]0,1[$, the numbers $\lambda_i/\mu$ for $i<n$ are nonnegative, and they add to $1$. The point
\[
y=\sum_{i=1}^{n-1}\frac{\lambda_i}{\mu}\,x_i
\]
is a convex combination of points of the interval $I$, so $y\in I$. The weighted average of all $n$ points is $\mu y+\lambda_nx_n$. Convexity at those two points gives
\[
f(\mu y+\lambda_nx_n)\le\mu f(y)+\lambda_nf(x_n).
\]
The induction hypothesis gives $f(y)\le\sum_{i<n}(\lambda_i/\mu)f(x_i)$. Multiply by the positive number $\mu$ and add $\lambda_nf(x_n)$. The right side becomes $\sum_{i=1}^n\lambda_if(x_i)$.
\end{proof}

When the bending is strict, equality forces every point to be the same point.

If $f''>0$ and every weight $\lambda_i$ is positive, equality in Jensen holds only when $x_1=\dots=x_n$. Let $m=\sum\lambda_ix_i$. Proposition~\ref{prop:convex}(b) gives $f(x_i)\ge f(m)+f'(m)(x_i-m)$, strictly unless $x_i=m$. Multiply by $\lambda_i$ and add. The right side collapses:
\[
\sum_{i=1}^n\lambda_if(x_i)\ge f(m)+f'(m)\sum_{i=1}^n\lambda_i(x_i-m)=f(m),
\]
which is Jensen again. Equality requires equality in every tangent estimate. Since each $\lambda_i$ is positive, every $x_i$ equals $m$.

The most used case is the comparison of the arithmetic and geometric means. The logarithm is strictly concave, so it lies under its tangents, and the tangent at the arithmetic mean is high enough to beat the average of the logarithms.

\begin{theorem}[AM--GM]\label{thm:amgm}
Let $x_1,\dots,x_n>0$ and let $\lambda_i>0$ with $\sum_{i=1}^n\lambda_i=1$. Then
\[
\prod_{i=1}^nx_i^{\lambda_i}\le\sum_{i=1}^n\lambda_ix_i,
\]
and in particular
\[
\sqrt[n]{x_1\cdots x_n}\le\frac{x_1+\dots+x_n}n.
\]
Equality holds if and only if $x_1=\dots=x_n$.
\end{theorem}
\begin{proof}
Let $m=\sum\lambda_ix_i$. Then $m>0$. The second derivative of $\ln$ is $-1/x^2<0$, so $-\ln$ has positive second derivative. Proposition~\ref{prop:convex}(b) applied to $-\ln$, and then multiplied by $-1$, gives the tangent bound for the logarithm:
\[
\ln x\le\ln m+\frac{x-m}m,
\]
with equality only at $x=m$. Apply this at each $x_i$, multiply by $\lambda_i$, and add:
\[
\sum_{i=1}^n\lambda_i\ln x_i\le\ln m+\frac1m\left(\sum_{i=1}^n\lambda_ix_i-m\right)=\ln m,
\]
because $\sum\lambda_ix_i=m$. The left side is the logarithm of $\prod x_i^{\lambda_i}$. The exponential function is increasing, so the products satisfy the claimed inequality. Equality holds only when each tangent estimate is an equality, hence only when every $x_i$ equals $m$.

The unweighted form is the case $\lambda_i=1/n$. For a check, take $1$, $4$ and $4$. The geometric mean is $\sqrt[3]{16}=2\sqrt[3]{2}\approx2.52$, and the arithmetic mean is $3$.
\end{proof}

The unweighted inequality extends to $x_i\ge0$. If some $x_i$ is zero, the geometric mean is zero and the arithmetic mean is nonnegative. If every $x_i$ is positive, the theorem applies. This extension is the one Carleman's inequality needs later.

Figure~\ref{fig:convex} shows both comparisons at once. On the left, a convex graph sits under a chord and above a tangent. On the right, $\ln$ sits under its tangent at the arithmetic mean $m=1.6$ of $0.5$, $1.2$ and $3.1$. The geometric mean is about $1.23$, and $\ln1.23$ is the height you get by averaging the three logarithms.

\begin{figure}[ht]
\centering
\fig{ch04-convex}
\caption{Left: a convex graph in charcoal sits under the rust chord and above the ochre tangent. Right: the proof of AM--GM for $x=(0.5,1.2,3.1)$. The charcoal points $(x_i,\ln x_i)$ lie under the ochre tangent to $\ln$ at the arithmetic mean $m=1.6$. Averaging the tangent values returns $\ln m$. Averaging the logarithms returns the logarithm of the geometric mean, about $\ln1.23$.}
\label{fig:convex}
\end{figure}

Cauchy already proved the unweighted comparison in Note~II of the \emph{Cours d'analyse} of $1821$. His argument climbs from $n$ to $2n$, up to a power of two, and then fills the missing places with copies of the mean so as to come back down. Exercise~\ref{ex:cauchyinduction} rebuilds that climb as a step from $n$ down to $n-1$. The same note contains the sum form of the inequality in the next section. Jensen's paper of $1906$ takes Cauchy's comparison of means as its starting point and isolates the chord property that implies it.

Johan Jensen was not a professor. He joined the Copenhagen telephone company in $1881$, became head of its technical department in $1890$, and kept that post until $1924$. The mathematics was spare-time work. The paper is ``Sur les fonctions convexes et les in\'egalit\'es entre les valeurs moyennes'', \emph{Acta Mathematica} $30$ ($1906$), $175$--$193$.

\portrait{cauchy}{Augustin-Louis Cauchy (1789--1857). Illustration from \emph{Das neunzehnte Jahrhundert in Bildnissen}, Berlin, 1901. Public domain, via Wikimedia Commons.}

\begin{examplex}[the largest box]
A rectangular box has twelve edges, four of each length $x$, $y$ and $z$, and the twelve lengths add to $12L$. Then $x+y+z=3L$. For positive $x,y,z$, Theorem~\ref{thm:amgm} gives
\[
xyz\le\left(\frac{x+y+z}3\right)^3=L^3,
\]
with equality when $x=y=z=L$, the cube. The inequality finds the maximum and the box that attains it. Exercise~\ref{ex:heron} is the same idea for the area of a triangle of fixed perimeter.
\end{examplex}

\section{Two lists}\label{sec:twolists}

The tangent compared values of one function. Two lists need a square that sees both lists at once.

\begin{think}
In the plane, $|\vec a\cdot\vec b|=\|\vec a\|\,\|\vec b\|\,|\cos\theta|$, and $|\cos\theta|\le1$. The cosine there is the ordinary trigonometric one. In $\R^n$ the angle is read off the dot product only after a bound that does not mention angles. Can you prove $\bigl|\sum a_ib_i\bigr|\le\sqrt{\sum a_i^2}\,\sqrt{\sum b_i^2}$ using only that a sum of squares is nonnegative?
\end{think}

Figure~\ref{fig:align} is the plane case, with numbers small enough to check by hand. The rust arrow $(4,3)$ has length $5$. Its shadow on the horizontal axis has length $4$, which is shorter. The charcoal arrow is $(4,0)$, of length $4$. The dot product is $16$, and the product of the lengths is $20$.

\begin{figure}[ht]
\centering
\fig{ch04-align}
\caption{Cauchy--Schwarz in the plane, for $\vec a=(4,3)$ and $\vec b=(4,0)$. The rust arrow has length $5$. The ochre segment is the perpendicular from its tip to the axis, of length $3$, so the shadow of $\vec a$ has length $4$. Then $\vec a\cdot\vec b=16$ and $\|\vec a\|\,\|\vec b\|=20$. Equality would need the perpendicular leg to vanish, which means the two arrows lie on the same line.}
\label{fig:align}
\end{figure}

The proof below does not use the angle. It uses a quadratic polynomial that cannot take a negative value.

\begin{theorem}[Cauchy--Schwarz]\label{thm:cs}
For real numbers $a_1,\dots,a_n$ and $b_1,\dots,b_n$,
\[
\left(\sum_{i=1}^na_ib_i\right)^2\le\left(\sum_{i=1}^na_i^2\right)\left(\sum_{i=1}^nb_i^2\right).
\]
Equality holds if and only if the two lists are linearly dependent, including the case where either list is entirely zero.
\end{theorem}
\begin{proof}
Write $A=\sum a_i^2$, $B=\sum a_ib_i$ and $C=\sum b_i^2$. If $C=0$, then every $b_i$ is $0$, so $B=0$ and both sides of the inequality are $0$. The lists are dependent.

Now suppose $C>0$. For every real $t$,
\[
\sum_{i=1}^n(a_i+tb_i)^2=A+2Bt+Ct^2\ge0,
\]
because it is a sum of squares. The quadratic $Ct^2+2Bt+A$ opens upwards. If its discriminant $4B^2-4AC$ were positive, the quadratic would be negative between its two roots. So the discriminant is at most $0$:
\[
B^2-AC\le0,
\]
which is the inequality.

Equality means the quadratic touches zero. It then has a root $t$, and at that root every square $(a_i+tb_i)^2$ vanishes, so $a_i=-tb_i$ for every $i$. The two lists are dependent. Conversely, if $a_i=sb_i$ for every $i$, or $b_i=sa_i$ for every $i$, both sides equal $s^2(\sum b_i^2)^2$ or the same identity with the lists swapped, and equality holds.
\end{proof}

The same polynomial works for an inner product. In particular, if $f$ and $g$ are continuous on a closed bounded interval, then
\[
\left(\int fg\right)^2\le\left(\int f^2\right)\left(\int g^2\right).
\]
The quadratic $t\mapsto\int(f+tg)^2$ is nonnegative, so its discriminant is at most $0$. If $g$ is not identically zero, equality means that this quadratic touches zero, hence $\int(f+tg)^2=0$ for some $t$. A continuous nonnegative function with integral zero vanishes everywhere, so $f+tg=0$.

One rearrangement of the same inequality is used often enough to have a nickname. For real $x_i$ and positive $y_i$,
\[
\sum_{i=1}^n\frac{x_i^2}{y_i}\ge\frac{\bigl(\sum_{i=1}^nx_i\bigr)^2}{\sum_{i=1}^ny_i}.
\]
Apply Theorem~\ref{thm:cs} to the lists $x_i/\sqrt{y_i}$ and $\sqrt{y_i}$. This is Engel's form, also called Titu's lemma in problem books. Equality holds when the two lists in that application are proportional, that is, when $x_i/y_i$ is the same for every $i$.

\begin{examplex}
How large can $x+2y+3z$ be when $x^2+y^2+z^2=1$? Theorem~\ref{thm:cs} with the lists $(x,y,z)$ and $(1,2,3)$ gives
\[
(x+2y+3z)^2\le(1+4+9)(x^2+y^2+z^2)=14.
\]
The left side is at most $\sqrt{14}$. Equality holds when $(x,y,z)$ is a scalar multiple of $(1,2,3)$. The condition $x^2+y^2+z^2=1$ fixes the scalar: the positive multiple is $(1,2,3)/\sqrt{14}$, and it gives the value $\sqrt{14}$. The negative multiple gives $-\sqrt{14}$, the minimum. Squaring the inequality hides that distinction. The dot product itself was defined in Chapter~\ref{ch:linear}. With $\|a\|=1$ and $b\ne0$, the inequality says $a\cdot b\le\|b\|$, and the value $\|b\|$ is attained when $a$ points the same way as $b$.
\end{examplex}

Cauchy's Note~II of $1821$ already contains the comparison of a sum with the sum of squares, in the form $(\sum a_i)^2\le n\sum a_i^2$, and the general sum form carries his name. Viktor Bunyakovsky, who had studied in Paris and treated that sum form as known, published the integral form in $1859$, in the memoir \emph{Sur quelques in\'egalit\'es concernant les int\'egrales ordinaires et les int\'egrales aux diff\'erences finies}. Hermann Schwarz found the integral form again in $1885$, twenty-six years later, while studying a minimal surface. His argument is the quadratic polynomial above. In the Russian literature the integral inequality is Bunyakovsky's.

The passage from Cauchy--Schwarz to a product of unequal powers is Young's inequality, and then H\"older's. Figure~\ref{fig:young} is Young's inequality with the algebra removed: two regions cover a rectangle.

\begin{figure}[ht]
\centering
\fig{ch04-young}
\caption{Young's inequality as areas, for $p=3$ and $q=3/2$. The curve $y=x^{p-1}$ is the same curve as $x=y^{q-1}$, because $(p-1)(q-1)=1$. The charcoal region under the curve from $0$ to $a$ has area $a^p/p$. The rust region to the left of the curve, from $0$ to $b$, has area $b^q/q$. Together they cover the dashed rectangle of area $ab$. The cover is exact, with overlap of area zero, precisely when the corner $(a,b)$ lies on the curve.}
\label{fig:young}
\end{figure}

\begin{theorem}[Young]\label{thm:young}
Let $p,q>1$ satisfy $1/p+1/q=1$, and let $a,b\ge0$. Then
\[
ab\le\frac{a^p}p+\frac{b^q}q,
\]
with equality if and only if $b=a^{p-1}$.
\end{theorem}
\begin{proof}
Figure~\ref{fig:young} is the reason for the two terms on the right: they are the area under $y=x^{p-1}$ and the area to the left of the same curve. The relation $1/p+1/q=1$ rearranges to $(p-1)(q-1)=1$, so $y=x^{p-1}$ is the same curve as $x=y^{q-1}$, and those areas are $a^p/p$ and $b^q/q$. The algebra below is the check that their sum is at least $ab$.

If $a=0$, the claim is $0\le b^q/q$, with equality exactly when $b=0$. And $a^{p-1}=0$ in that case. The case $b=0$ is symmetric, and equality forces $a=0$.

Now let $a>0$, and set $\varphi(b)=a^p/p+b^q/q-ab$ for $b\ge0$. Then $\varphi'(b)=b^{q-1}-a$. The derivative vanishes when $b^{q-1}=a$. Since $q-1=1/(p-1)$, this value is $b=a^{p-1}$, and it is positive. The second derivative $\varphi''(b)=(q-1)b^{q-2}$ is positive for $b>0$, so $\varphi'$ is strictly increasing: negative before $a^{p-1}$ and positive after. Thus $\varphi$ has a strict minimum there. The minimum value is
\[
\varphi(a^{p-1})=\frac{a^p}p+\frac{a^{(p-1)q}}q-a\cdot a^{p-1}.
\]
The exponent $(p-1)q$ equals $p$, so $a^{(p-1)q}=a^p$, and the right side is $a^p(1/p+1/q-1)=0$. Therefore $\varphi(b)\ge0$, with equality only at $b=a^{p-1}$.
\end{proof}

Exercise~\ref{ex:holder} proves the same inequality from weighted AM--GM, with no picture. Either proof gives H\"older's inequality, by applying Young to normalised terms.

\begin{theorem}[H\"older]\label{thm:holder}
Let $p,q>1$ satisfy $1/p+1/q=1$, and let $a_i,b_i\ge0$. Then
\[
\sum_{i=1}^na_ib_i\le\left(\sum_{i=1}^na_i^p\right)^{1/p}\left(\sum_{i=1}^nb_i^q\right)^{1/q}.
\]
If both sides are positive, equality holds if and only if $a_i^p=c\,b_i^q$ for a constant $c\ge0$ and every $i$.
\end{theorem}
\begin{proof}
Let $A=(\sum a_i^p)^{1/p}$ and $B=(\sum b_i^q)^{1/q}$. If $A=0$, then every $a_i$ is $0$ and both sides vanish. If $B=0$, every $b_i$ is $0$ and both sides vanish. Now suppose $A>0$ and $B>0$. Theorem~\ref{thm:young} at the nonnegative numbers $a_i/A$ and $b_i/B$ gives
\[
\frac{a_i}{A}\cdot\frac{b_i}{B}\le\frac1p\left(\frac{a_i}{A}\right)^p+\frac1q\left(\frac{b_i}{B}\right)^q.
\]
Add over $i$. The left side is $(\sum a_ib_i)/(AB)$. On the right,
\[
\frac1p\frac{\sum a_i^p}{A^p}+\frac1q\frac{\sum b_i^q}{B^q}=\frac1p+\frac1q=1.
\]
Multiply by the positive number $AB$.

Equality, when $A$ and $B$ are positive, requires equality in Young's inequality at every $i$. That is $b_i/B=(a_i/A)^{p-1}$. Raise both sides to the power $q$. Since $(p-1)q=p$,
\[
\frac{b_i^q}{B^q}=\frac{a_i^p}{A^p},
\]
so $a_i^p$ is the constant multiple $(A/B)^p$ of $b_i^q$.
\end{proof}

Cauchy--Schwarz is the case $p=q=2$, for nonnegative lists. The signed equality case, linear dependence with a negative scalar, is the one in Theorem~\ref{thm:cs}. H\"older as stated here does not see signs, because the powers $a_i^p$ for a non-integer $p$ are taken on nonnegative numbers.

\begin{think}
You are paid for three jobs of lengths $a_1\le a_2\le a_3$ at three hourly rates $b_1\le b_2\le b_3$, and you choose which rate goes with which job. Which matching pays the most, and which pays the least? One swap of two jobs is enough to see the direction.
\end{think}

Paying the long job at the high rate is the expensive way to staff the week. Paying it at the low rate is the cheap way. Every other matching sits in between. One swap of a mismatched pair never decreases the total, and that single comparison is the induction.

\begin{proposition}[rearrangement]\label{prop:rearr}
Let $a_1\le\dots\le a_n$ and $b_1\le\dots\le b_n$. Then for every permutation $\sigma$,
\[
\sum_{i=1}^na_ib_{n+1-i}\le\sum_{i=1}^na_ib_{\sigma(i)}\le\sum_{i=1}^na_ib_i.
\]
The sum of products is largest when the two lists are sorted the same way, and smallest when they are sorted in opposite ways.
\end{proposition}
\begin{proof}
The upper bound is an induction on $n$. For $n=1$ there is nothing to compare. Let $k$ be the index with $\sigma(k)=n$, so $b_n$ is paired with $a_k$. If $k=n$, the permutation already fixes $n$, and the induction hypothesis applies to the first $n-1$ places.

If $k\ne n$, exchange the values of $\sigma$ at $k$ and at $n$. The sum changes by
\begin{align*}
a_kb_{\sigma(n)}+a_nb_n-a_kb_n-a_nb_{\sigma(n)}
&=(a_n-a_k)(b_n-b_{\sigma(n)}).
\end{align*}
Here $a_n\ge a_k$. Also $b_n$ is the largest $b$-value, so $b_n\ge b_{\sigma(n)}$. The product is nonnegative: the swap does not decrease the sum. The new permutation fixes $n$. On $\{1,\dots,n-1\}$ it is some permutation, and the induction hypothesis says that pairing is at most $\sum_{i<n}a_ib_i$. Adding $a_nb_n$ finishes the upper bound.

The lower bound is the upper bound applied to the lists $a_i$ and $-b_{n+1-i}$. Both of those lists are sorted in increasing order. The largest sum for that second list is the smallest sum for $b$.
\end{proof}

Figure~\ref{fig:rearrange} checks the two extremes on $(1,2,3)$ and $(1,2,4)$. The same-order sum is $1+4+12=17$. The opposite-order sum is $4+4+3=11$. Any other matching, for instance $1\cdot2+2\cdot1+3\cdot4=16$, lands between them.

\begin{figure}[ht]
\centering
\fig{ch04-rearrange}
\caption{Rearrangement for the lists $(1,2,3)$ and $(1,2,4)$. Pairing in the same order gives the products $1$, $4$ and $12$, which add to $17$. Pairing in opposite orders gives $4$, $4$ and $3$, which add to $11$. Proposition~\ref{prop:rearr} says every other matching lands between these two sums.}
\label{fig:rearrange}
\end{figure}

\begin{examplex}
Let $x,y,z>0$. Sort them as the first list, and sort $1/x$, $1/y$, $1/z$ as the second. The second list runs in the opposite order. The opposite pairing includes the sum of each number with its own reciprocal,
\[
x\cdot\frac1x+y\cdot\frac1y+z\cdot\frac1z=3,
\]
and Proposition~\ref{prop:rearr} says this is the smallest possible sum of products. Every other pairing is at least $3$. One of them is
\[
\frac xy+\frac yz+\frac zx\ge3.
\]
The hypothesis $xyz=1$ is not used. Replacing $(x,y,z)$ by $(a^2b,b^2c,c^2a)$ gives the same bound in the form $a^2/(bc)+b^2/(ca)+c^2/(ab)\ge3$, for positive $a,b,c$.
\end{examplex}

\section{Power means and norms}\label{sec:means}

H\"older compared a product of two lists. The power means compare one list with itself, at different exponents. They line up.

\begin{theorem}[power means]\label{thm:powermeans}
Let $x_1,\dots,x_n>0$ and let $\lambda_i>0$ with $\sum_{i=1}^n\lambda_i=1$. For $p\ne0$ set
\[
M_p=\left(\sum_{i=1}^n\lambda_ix_i^p\right)^{1/p},\qquad M_0=\prod_{i=1}^nx_i^{\lambda_i}.
\]
The map $p\mapsto M_p$ is non-decreasing on $\R$, and it is strictly increasing unless $x_1=\dots=x_n$. Also $M_p\to M_0$ as $p\to0$, $M_p\to\max_ix_i$ as $p\to+\infty$, and $M_p\to\min_ix_i$ as $p\to-\infty$.
\end{theorem}
\begin{proof}
\emph{Positive exponents.} Let $0<p<q$. The map $t\mapsto t^{q/p}$ on $]0,\infty[$ has second derivative $(q/p)(q/p-1)t^{q/p-2}>0$, so it is strictly convex. Jensen's inequality at the points $x_i^p$, with these weights, gives
\[
\left(\sum_{i=1}^n\lambda_ix_i^p\right)^{q/p}\le\sum_{i=1}^n\lambda_i(x_i^p)^{q/p}=\sum_{i=1}^n\lambda_ix_i^q.
\]
That is $M_p^q\le M_q^q$. Both means are positive, so $M_p\le M_q$. Equality in Jensen for a strictly convex function and positive weights forces the points $x_i^p$ to be equal, hence the $x_i$ to be equal, because $t\mapsto t^p$ is one-to-one on $]0,\infty[$.

\emph{Across zero, from AM--GM.} Let $q>0$. Apply Theorem~\ref{thm:amgm} to the positive numbers $x_i^q$:
\[
\prod_{i=1}^n(x_i^q)^{\lambda_i}\le\sum_{i=1}^n\lambda_ix_i^q.
\]
The left side is $M_0^q$ and the right side is $M_q^q$, so $M_0\le M_q$, with equality only when all the $x_i$ are equal.

\emph{Negative exponents, by reciprocals.} For every real $r$, including $0$,
\[
M_r(x_1,\dots,x_n)=\frac1{M_{-r}(1/x_1,\dots,1/x_n)}.
\]
Indeed, if $r\ne0$,
\[
M_{-r}(1/x)=\left(\sum_{i=1}^n\lambda_ix_i^{r}\right)^{-1/r}=\frac1{M_r(x)},
\]
and the case $r=0$ is $\prod(1/x_i)^{\lambda_i}=1/M_0(x)$. Now let $p<q<0$. Then $0<-q<-p$, so the positive-exponent case gives $M_{-q}(1/x)\le M_{-p}(1/x)$. All of these means are positive. Taking reciprocals reverses the inequality:
\[
M_q(x)=\frac1{M_{-q}(1/x)}\ge\frac1{M_{-p}(1/x)}=M_p(x).
\]
The comparison $M_p\le M_0$ for $p<0$ is the same trick applied to $M_0(1/x)\le M_{-p}(1/x)$. A strict inequality between positive numbers stays strict when you take reciprocals, so the means are strictly ordered unless all the $x_i$ are equal.

\emph{The limit at infinity.} Let $\bar x=\max_ix_i$, attained at an index $i_0$. For $p>0$,
\[
\lambda_{i_0}\bar x^p\le\sum_{i=1}^n\lambda_ix_i^p\le\bar x^p,
\]
because every $x_i\le\bar x$ and the weights add to $1$. Take $p$th roots:
\[
\lambda_{i_0}^{1/p}\,\bar x\le M_p\le\bar x.
\]
As $p\to+\infty$, the factor $\lambda_{i_0}^{1/p}$ tends to $1$. Squeeze the two ends: $M_p\to\bar x$. The limit as $p\to-\infty$ is the reciprocal identity applied to $1/x_i$, whose maximum is $1/\min x_i$.

\emph{The limit at zero.} Set $t_i=\ln x_i$ and $\varphi(p)=\ln\sum_{i=1}^n\lambda_ie^{pt_i}$, so that $\ln M_p=\varphi(p)/p$ for $p\ne0$. Then $\varphi(0)=\ln\sum\lambda_i=\ln1=0$. Differentiating under the sum, the derivative of $\sum\lambda_ie^{pt_i}$ at $p=0$ is $\sum\lambda_it_i$, and the derivative of $\ln$ at $1$ is $1$, so $\varphi'(0)=\sum\lambda_it_i$. Therefore
\[
\frac{\varphi(p)}{p}=\frac{\varphi(p)-\varphi(0)}{p-0}\to\varphi'(0)=\sum_{i=1}^n\lambda_i\ln x_i=\ln M_0
\]
as $p\to0$. Exponentiate.
\end{proof}

\begin{figure}[ht]
\centering
\fig{ch04-powermeans}
\caption{Power means of $(1,2,4,8)$ with equal weights. The charcoal curve rises from the minimum $1$ to the maximum $8$. The rust dots are the harmonic mean $H=M_{-1}\approx2.13$, the geometric mean $G=M_0\approx2.83$, the arithmetic mean $A=M_1=3.75$ and the quadratic mean $Q=M_2\approx4.61$. The chain $H\le G\le A\le Q$ is the part of Theorem~\ref{thm:powermeans} used every day.}
\label{fig:powermeans}
\end{figure}

For a check with equal weights on $(1,2,4,8)$: the harmonic mean is $4/(1+1/2+1/4+1/8)=4/(15/8)=32/15\approx2.13$, the geometric mean is $(64)^{1/4}=\sqrt8\approx2.83$, the arithmetic mean is $15/4=3.75$, and the quadratic mean is $\sqrt{(1+4+16+64)/4}=\sqrt{85/4}\approx4.61$. Figure~\ref{fig:powermeans} is that list, drawn for every real exponent.

Chebyshev's sum inequality is a different comparison, for two lists that rise together. The next chapter gives Chebyshev's name to an inequality about random variables. The one here is about sorted numbers.

\begin{proposition}[Chebyshev's sum inequality]\label{prop:chebsum}
If $a_1\le\dots\le a_n$ and $b_1\le\dots\le b_n$, then
\[
\frac1n\sum_{i=1}^na_ib_i\ge\left(\frac1n\sum_{i=1}^na_i\right)\left(\frac1n\sum_{i=1}^nb_i\right).
\]
If one list is increasing and the other is decreasing, the inequality reverses.
\end{proposition}
\begin{proof}
For every pair of indices $i$ and $j$, the differences $a_i-a_j$ and $b_i-b_j$ have the same sign, so $(a_i-a_j)(b_i-b_j)\ge0$. Expand and add over all $i$ and $j$, including the diagonal where the product is zero:
\begin{align*}
\sum_{i=1}^n\sum_{j=1}^n(a_i-a_j)(b_i-b_j)
&=\sum_{i,j}\bigl(a_ib_i-a_ib_j-a_jb_i+a_jb_j\bigr)\\
&=2n\sum_{i=1}^na_ib_i-2\left(\sum_{i=1}^na_i\right)\left(\sum_{j=1}^nb_j\right).
\end{align*}
The double sum is nonnegative. Divide by $2n^2$. If the lists are sorted in opposite directions, every product $(a_i-a_j)(b_i-b_j)$ is nonpositive, and the inequality reverses.
\end{proof}

In words: when the lists rise together, the mean of the products is at least the product of the means. That gap is a covariance. Volume~4 defines $\operatorname{Cov}(X,Y)$ as $\mathbb E[XY]-\mathbb E[X]\mathbb E[Y]$, and the same sign computation shows the covariance is nonnegative when $X$ and $Y$ are increasing functions of one variable. Exercise~\ref{ex:chebsumapp} uses the proposition to compare power means of equal weights, with no appeal to Jensen.

The last inequality of this section is the triangle inequality for the size
\[
\|a\|_p=\left(\sum_{i=1}^n|a_i|^p\right)^{1/p},\qquad p\ge1.
\]
For $p=2$ this is the Euclidean length from Chapter~\ref{ch:linear}. Hermann Minkowski proved the general case in \emph{Geometrie der Zahlen} (Teubner, Leipzig, $1896$), where the same expression appears as a way of measuring a convex body used in the geometry of numbers.

\portrait{minkowski}{Hermann Minkowski (1864--1909). Photograph from \emph{Raum und Zeit}, 1909. Public domain, via Wikimedia Commons.}

\begin{theorem}[Minkowski]\label{thm:minkowski}
For $p\ge1$ and $a_i,b_i\ge0$,
\[
\left(\sum_{i=1}^n(a_i+b_i)^p\right)^{1/p}\le\left(\sum_{i=1}^na_i^p\right)^{1/p}+\left(\sum_{i=1}^nb_i^p\right)^{1/p}.
\]
For $p=1$ the two sides are equal. For $p>1$, equality holds if and only if the two lists are proportional, including the case where either list is entirely zero.
\end{theorem}
\begin{proof}
For $p=1$ the identity $(a_i+b_i)=a_i+b_i$ sums to an equality. Now let $p>1$, and set $q=p/(p-1)$, so that $1/p+1/q=1$ and $(p-1)q=p$. Let $S=\sum_{i=1}^n(a_i+b_i)^p$. If $S=0$, then every $a_i$ and $b_i$ is $0$ and both sides vanish.

If $S>0$, split each term:
\[
(a_i+b_i)^p=a_i(a_i+b_i)^{p-1}+b_i(a_i+b_i)^{p-1}.
\]
Apply Theorem~\ref{thm:holder} to each of the two resulting sums, with the lists $a_i$ and $(a_i+b_i)^{p-1}$, and with $b_i$ and $(a_i+b_i)^{p-1}$:
\begin{align*}
S&\le\left(\sum_{i=1}^na_i^p\right)^{1/p}\left(\sum_{i=1}^n(a_i+b_i)^{(p-1)q}\right)^{1/q}\\
&\qquad+\left(\sum_{i=1}^nb_i^p\right)^{1/p}\left(\sum_{i=1}^n(a_i+b_i)^{(p-1)q}\right)^{1/q}.
\end{align*}
The exponent $(p-1)q$ equals $p$, so each second factor is $S^{1/q}$. Divide by this positive number. Since $1-1/q=1/p$,
\[
S^{1/p}\le\left(\sum_{i=1}^na_i^p\right)^{1/p}+\left(\sum_{i=1}^nb_i^p\right)^{1/p}.
\]

For the equality case when $p>1$ and $S>0$, equality must hold in both applications of H\"older. Theorem~\ref{thm:holder}, applied to $a_i$ and $(a_i+b_i)^{p-1}$, says that $a_i^p=c(a_i+b_i)^p$ for some constant $c\ge0$ and every $i$. Taking $p$th roots gives $a_i=c^{1/p}(a_i+b_i)$. If $c^{1/p}\ne1$, then $b_i$ is a constant multiple of $a_i$. If $c^{1/p}=1$, then $b_i=0$ for every $i$. If $c=0$, then $a_i=0$ for every $i$. In all three cases the lists are proportional. Conversely, if one list is a nonnegative multiple of the other, both sides are equal: if $b_i=\lambda a_i$ for some $\lambda\ge0$, both sides equal $(1+\lambda)(\sum a_i^p)^{1/p}$, and the case of a zero first list is the same identity with the lists swapped.
\end{proof}

For arbitrary real lists $u$ and $v$, apply the theorem to $|u_i|$ and $|v_i|$, and use $|u_i+v_i|\le|u_i|+|v_i|$ in each term. The result is $\|u+v\|_p\le\|u\|_p+\|v\|_p$. At $p=2$ this says that a straight line is a shortest path in $\R^n$. As $p\to\infty$, Theorem~\ref{thm:powermeans} says that the $p$-norm of a fixed list tends to $\max_i|a_i|$, and the triangle inequality for the maximum is the observation that $|u_i+v_i|\le\max|u|+\max|v|$ for every $i$. Volume~7 builds the spaces $L^p$ on H\"older's inequality and on this one.

\section{Carleman's constant}\label{sec:carleman}

Every constant so far came from a square or a tangent, and the equality case said when the constant could not be improved. The constant $e$ is of that kind. It enters through the sequence $(1+1/k)^k$, which is strictly less than $e$, and a single family of sequences forces every smaller constant to fail.

Let $a_n\ge0$ and suppose $\sum_{n=1}^\infty a_n$ converges. Carleman's inequality, proved in Exercise~\ref{ex:carleman}, says
\[
\sum_{n=1}^\infty(a_1a_2\cdots a_n)^{1/n}\le e\sum_{n=1}^\infty a_n.
\]
Torsten Carleman stated it in the proceedings of the Scandinavian congress, published at Helsinki in $1923$, in the course of work on quasi-analytic functions. George P\'olya's proof, the one the exercise follows, was recorded in the \emph{Proceedings of the London Mathematical Society} $(2)$ $24$ ($1926$), page~lvii. It is weighted AM--GM with a carefully chosen weight.

A first attempt shows why a weight is needed. AM--GM on the first $n$ terms gives $(a_1\cdots a_n)^{1/n}\le(1/n)\sum_{k\le n}a_k$. Add over $n$. The coefficient of a fixed $a_k$ is $\sum_{n\ge k}1/n$, and the harmonic series diverges. The resulting upper bound is infinite. It does not prove convergence of anything.

Put a weight inside the geometric mean, and divide it out afterwards. Set
\[
c_k=\frac{(k+1)^k}{k^{k-1}}.
\]
The product telescopes. For a check at $n=3$: $c_1=2$, $c_2=9/2$ and $c_3=64/9$, and $2\cdot(9/2)\cdot(64/9)=64=4^3$. In general the intermediate powers cancel:
\[
\prod_{k=1}^nc_k=\frac{2^1\cdot3^2\cdots(n+1)^n}{1^0\cdot2^1\cdots n^{n-1}}=(n+1)^n.
\]
Weighted AM--GM on the $n$ positive numbers $c_ka_k$, or the nonnegative extension if some $a_k$ vanish, gives
\[
(a_1\cdots a_n)^{1/n}=\frac{(c_1a_1\cdots c_na_n)^{1/n}}{n+1}\le\frac1{n(n+1)}\sum_{k=1}^nc_ka_k.
\]
Add over $n$. Every term is nonnegative, so Proposition~\ref{prop:safe-sums} allows the order of summation to be swapped. The coefficient of $c_ka_k$ is
\[
\sum_{n\ge k}\frac1{n(n+1)}=\sum_{n\ge k}\left(\frac1n-\frac1{n+1}\right)=\frac1k,
\]
a telescoping sum. Therefore the left side of Carleman's inequality is at most $\sum_k(c_k/k)a_k$. And
\[
\frac{c_k}k=\left(1+\frac1k\right)^k.
\]
The tangent to $\ln$ at $1$ lies above the graph: for $u>0$, $\ln(1+u)<u$. At $u=1/k$,
\[
k\ln\left(1+\frac1k\right)<1,\qquad\left(1+\frac1k\right)^k<e.
\]
Each coefficient is strictly less than $e$, and the series of geometric means is at most $e$ times $\sum a_k$. Exercise~\ref{ex:carleman} asks you to write this chain in order. Exercise~\ref{ex:euler} shows that $(1+1/n)^n$ is strictly increasing, which with the upper bound $e$ is one way to see that the limit exists.

The constant is sharp, but the sequence that shows it is slow. Take $a_n=1/n$ for $n\le N$ and $a_n=0$ afterwards. The left side divided by $\sum_{n\le N}a_n$ is about $1.52$ when $N=10$, about $1.94$ when $N=100$, and about $2.29$ when $N=10^4$. The first challenge proves that this quotient tends to $e$. A smaller constant in front of $\sum a_n$ fails for all large $N$. The numerical values alone climb too slowly to make that conclusion; the challenge uses Stirling's formula from Theorem~\ref{thm:stirlingC}.

For a sequence that decays faster, the inequality has more room. The first eight geometric means of $1,1/4,1/9,\dots,1/64$ add to about $2.42$. The number $e$ times the sum of those eight terms $1/n^2$ is about $4.15$.

\section{Log-convexity of the Gamma function}\label{sec:gamma-convexity}

The same H\"older estimate, written for integrals, singles out the Gamma function among the functions that interpolate the factorials.

Chapter~\ref{ch:calculus} defined $\Gamma(x)=\int_0^\infty t^{x-1}e^{-t}\,\dd t$ for $x>0$ and proved $\Gamma(x+1)=x\Gamma(x)$, so $\Gamma(n+1)=n!$. Many other functions satisfy the same recurrence. An inequality removes the ambiguity.

\begin{proposition}[log-convexity of Gamma]\label{prop:gamma-convexity}
For $x,y>0$ and $0\le\lambda\le1$,
\[
\Gamma\bigl(\lambda x+(1-\lambda)y\bigr)\le\Gamma(x)^\lambda\,\Gamma(y)^{1-\lambda}.
\]
Equivalently, $\ln\Gamma$ is convex on $]0,\infty[$.
\end{proposition}
\begin{proof}
The endpoint cases $\lambda=0$ and $\lambda=1$ are equalities. Fix $0<\lambda<1$, and set $p=1/\lambda$ and $q=1/(1-\lambda)$. Then $p,q>1$ and $1/p+1/q=1$.

First apply Young on a finite interval, where the integrals are ordinary. Let $f$ and $g$ be continuous on $[\varepsilon,A]$ with $0<\varepsilon<A$, and set $U=(\int_\varepsilon^A|f|^p)^{1/p}$ and $V=(\int_\varepsilon^A|g|^q)^{1/q}$. If $U=0$, then the continuous nonnegative function $|f|^p$ has integral zero, so $f=0$ and $\int|fg|=0$. The case $V=0$ is the same. If $U>0$ and $V>0$, Theorem~\ref{thm:young} at the numbers $|f(t)|/U$ and $|g(t)|/V$, integrated over $[\varepsilon,A]$, gives
\[
\frac{\int_\varepsilon^A|fg|}{UV}\le\frac1p\frac{\int_\varepsilon^A|f|^p}{U^p}+\frac1q\frac{\int_\varepsilon^A|g|^q}{V^q}=1.
\]

Now choose the integrands that rebuild $\Gamma$. Set
\[
f(t)=\bigl(t^{x-1}e^{-t}\bigr)^\lambda,\qquad g(t)=\bigl(t^{y-1}e^{-t}\bigr)^{1-\lambda}.
\]
Their product is $t^{\lambda x+(1-\lambda)y-1}e^{-t}$, the integrand of $\Gamma$ at $\lambda x+(1-\lambda)y$. Also $|f|^p=t^{x-1}e^{-t}$ and $|g|^q=t^{y-1}e^{-t}$, because $\lambda p=1$ and $(1-\lambda)q=1$. Let $\varepsilon\to0^+$ and $A\to\infty$. All three integrals increase to the corresponding integrals from $0$ to $\infty$, which converge by Proposition~\ref{prop:gammafn}. The inequality passes to the limit:
\[
\Gamma\bigl(\lambda x+(1-\lambda)y\bigr)\le\Gamma(x)^{1/p}\,\Gamma(y)^{1/q}=\Gamma(x)^\lambda\,\Gamma(y)^{1-\lambda}.
\]
The logarithm is increasing, so $\ln\Gamma$ satisfies the chord inequality. That is the convexity of $\ln\Gamma$.
\end{proof}

The shape of the proof is worth keeping. An average in the argument of $\Gamma$ becomes a product in the integrand, because the power $t^{\lambda x+(1-\lambda)y}$ factors. H\"older is built to separate such products.

Bohr and Mollerup proved in $1922$ that this property characterises $\Gamma$. Among positive functions on $]0,\infty[$ with $f(1)=1$ and $f(x+1)=xf(x)$, log-convexity leaves only $\Gamma$. Chapter~\ref{ch:calculus} records that fact as an outlook. Volume~5 proves it.

\section{Workshop: the cost of being unequal}\label{sec:slack}

An inequality that names the maximum is more useful when it also prices the distance from that maximum.

\begin{examplex}[the bill for an unequal rectangle]
Two nonnegative side lengths $x$ and $y$ have a fixed sum $s$. Their rectangle has area
\[
xy=\frac{s^2}4-\frac{(x-y)^2}4.
\]
The square of side $s/2$ has area $s^2/4$. An imbalance of $d=|x-y|$ costs exactly $d^2/4$ in area. If $s=10$ and the area is at least $24$, then $(x-y)^2\le4$, so the sides differ by at most $2$. A near-maximal area forces near-equal sides. Both sides of the identity are areas, lengths squared.
\end{examplex}

\begin{figure}[ht]
\centering
\fig{review04-slack}
\caption{For side-sum $10$, the area is $x(10-x)$ on $0\le x\le10$. The maximum $25$ is at $x=5$. The dashed line at height $24$ meets the curve at $x=4$ and $x=6$. Between those roots the area is at least $24$, so the sides differ by at most $2$.}
\label{fig:slack}
\end{figure}

The same receipt works for a product of a sum and a sum of reciprocals. For $x>0$ and $y>0$,
\[
(x+y)\left(\frac1x+\frac1y\right)-4=\frac{(x-y)^2}{xy}.
\]
The right side is nonnegative, so the left side is at least $4$, with equality exactly when $x=y$. The identity is exact in the other direction as well: the left side is at most $4+\eta$, for a given $\eta\ge0$, if and only if $|x-y|\le\sqrt{\eta\,xy}$. The bound controls a relative gap. It does not bound $|x-y|$ by a fixed length until a scale for $xy$ is supplied.

Check the algebra once. The common denominator $xy$ is positive, and
\[
(x+y)\left(\frac1x+\frac1y\right)-4=\frac{(x+y)^2-4xy}{xy}=\frac{(x-y)^2}{xy}.
\]

\begin{reminder}[the zero list in Cauchy--Schwarz]
Equality in Theorem~\ref{thm:cs} means the two lists are linearly dependent. That includes the case where either list is zero. If the second list is known to be nonzero, the first is a scalar multiple of it. For the maximum of $a\cdot b$ with $\|a\|=1$ and $b\ne0$, only the positive multiple $a=b/\|b\|$ attains the maximum. The negative multiple attains the minimum. Squaring both sides of an inequality can hide which one you found.
\end{reminder}

\section*{What this chapter prepares}

A square gave the comparison of two means. A chord gave Jensen's inequality, and the tangent to the logarithm gave the comparison for any number of positive terms. The same square, written as a quadratic in an auxiliary variable $t$, gave Cauchy--Schwarz. Young and H\"older extended it to unequal powers, and Minkowski's inequality is the triangle inequality for the $p$-norm when $p\ge1$. The constant $e$ is what those estimates leave in front of a series of geometric means. Equality cases are part of each statement. They are how one knows that a constant cannot be improved, as with the $e$ in Exercise~\ref{ex:carleman}.

The next chapter puts nonnegative weights on a finite set and calls them probabilities. Jensen's inequality becomes a statement about an average against those weights. The first limit in that chapter is a count from Chapter~\ref{ch:counting}: if $n$ guests collect their hats at random, the probability that nobody receives their own hat tends to $1/e$.

Volume~4 reads Jensen as $f(\mathbb E[X])\le\mathbb E[f(X)]$ and uses Cauchy--Schwarz on covariances. Volume~5 estimates series and integrals by the same tangents. Volume~7 builds $L^p$ from H\"older and Minkowski. Volume~11 turns estimates of this kind into concentration. The primes, in Volumes~6 and~14, are where the same comparisons bound sums that no longer have a closed form.

\section{Exercises}

\begin{exercise}[$\star$]\label{ex:8abc}
Show that $(a+b)(b+c)(c+a)\ge8abc$ for $a,b,c>0$.
\end{exercise}

\begin{exercise}[$\star$]\label{ex:sqrtk}
Show that $\displaystyle\sum_{k=1}^n\sqrt k\le n\sqrt{\frac{n+1}2}$.
\end{exercise}

\begin{exercise}[$\star$]\label{ex:intcs}
Let $f$ be continuous and positive on $[0,1]$. Show that $\displaystyle\int_0^1f\cdot\int_0^1\frac1f\ge1$. When is there equality?
\end{exercise}

\begin{exercise}[$\star\star$]\label{ex:euler}
Show that the sequence $\left(1+\frac1n\right)^n$ is strictly increasing.
\end{exercise}

\begin{exercise}[$\star\star$]\label{ex:nesbitt}
\emph{Nesbitt's inequality.} Show that $\dfrac a{b+c}+\dfrac b{c+a}+\dfrac c{a+b}\ge\dfrac32$ for $a,b,c>0$.
\end{exercise}

\begin{exercise}[$\star\star$]\label{ex:triangle}
Show that in every triangle with angles $A,B,C$, $\sin A+\sin B+\sin C\le\dfrac{3\sqrt3}2$.
\end{exercise}

\begin{exercise}[$\star\star$]\label{ex:aabbcc}
Show that $a^ab^bc^c\ge(abc)^{(a+b+c)/3}$ for $a,b,c>0$.
\end{exercise}

\begin{exercise}[$\star\star$]\label{ex:holder}
\begin{enumerate}[label=(\alph*)]
\item \emph{Young, from AM--GM.} Let $p,q>1$ with $\frac1p+\frac1q=1$. Prove Theorem~\ref{thm:young} again, this time from weighted AM--GM, with no areas.
\item \emph{H\"older.} Deduce Theorem~\ref{thm:holder}.
\end{enumerate}
\end{exercise}

\begin{exercise}[$\star\star\star$]\label{ex:carleman}
\emph{Carleman's inequality.} Let $(a_n)_{n\ge1}$ be non-negative with $\sum a_n<\infty$. Show that
\[
\sum_{n=1}^\infty\left(a_1a_2\cdots a_n\right)^{1/n}\le e\sum_{n=1}^\infty a_n.
\]
\end{exercise}

\begin{beauty}
Exercise~\ref{ex:carleman} is the opening question. The weights produce $\bigl(1+\frac1k\bigr)^k$, which is strictly less than $e$ by the tangent to $\ln$, and which Exercise~\ref{ex:euler} shows is increasing. The first challenge proves that the constant cannot be replaced by anything smaller. In Volume~4, Jensen reads $f(\mathbb E[X])\le\mathbb E[f(X)]$, and Cauchy--Schwarz bounds a covariance. In Volume~7, H\"older is the comparison that builds $L^p$.
\end{beauty}

\subsection*{Further exercises}

\begin{exercise}[$\star$]\label{ex:means3}
For $a,b,c>0$ show directly, without Theorem~\ref{thm:powermeans}, that
\[
\frac3{\frac1a+\frac1b+\frac1c}\le\sqrt[3]{abc}\le\frac{a+b+c}3\le\sqrt{\frac{a^2+b^2+c^2}3}.
\]
\end{exercise}

\begin{exercise}[$\star$]\label{ex:bernoulli}
\begin{enumerate}[label=(\alph*)]
\item Use Bernoulli's inequality to show that $\bigl(1+\frac1n\bigr)^n\ge2$ and $\bigl(1-\frac1{n^2}\bigr)^n\ge1-\frac1n$ for $n\ge1$.
\item Show that $\sqrt[n]n\to1$. \emph{Hint: write $\sqrt[n]n=1+x_n$ and keep two terms of the binomial expansion of $(1+x_n)^n$.}
\end{enumerate}
\end{exercise}

\begin{exercise}[$\star\star$]\label{ex:lagrangeid}
\emph{Lagrange's identity.} Show that for real $a_i,b_i$,
\[
\Bigl(\sum_{i=1}^na_i^2\Bigr)\Bigl(\sum_{i=1}^nb_i^2\Bigr)-\Bigl(\sum_{i=1}^na_ib_i\Bigr)^2=\sum_{1\le i<j\le n}(a_ib_j-a_jb_i)^2,
\]
and deduce Cauchy--Schwarz with its equality case. What does the identity say for $n=2$ about complex numbers?
\end{exercise}

\begin{exercise}[$\star\star$]\label{ex:cauchyinduction}
\emph{Cauchy's proof of AM--GM (1821).} Let $G_n\le A_n$ denote the inequality $\sqrt[n]{x_1\cdots x_n}\le\frac{x_1+\dots+x_n}n$ for positive numbers.
\begin{enumerate}[label=(\alph*)]
\item Prove $G_2\le A_2$, and show that $G_n\le A_n$ implies $G_{2n}\le A_{2n}$.
\item Show that $G_n\le A_n$ implies $G_{n-1}\le A_{n-1}$, by applying it to $x_1,\dots,x_{n-1}$ and $x_n=\frac{x_1+\dots+x_{n-1}}{n-1}$.
\item Conclude that $G_n\le A_n$ holds for every $n$.
\end{enumerate}
\end{exercise}

\begin{exercise}[$\star\star$]\label{ex:heron}
Heron's formula gives the area of a triangle with sides $a,b,c$ and half-perimeter $s=\frac{a+b+c}2$ as $\sqrt{s(s-a)(s-b)(s-c)}$. Show that among all
triangles of perimeter $2s$, the equilateral one has the largest area, and compute that area.
\end{exercise}

\begin{exercise}[$\star\star$]\label{ex:entropy}
\emph{Bridge: Convexity and finite probability weights.}
Let $p_1,\dots,p_n>0$ with $\sum p_i=1$. Show that
\[
H=\sum_{i=1}^np_i\ln\frac1{p_i}\le\ln n,
\]
with equality only for $p_i=\frac1n$. (In Volume~4, $H$ is the entropy of a probability distribution: the uniform distribution is the most unpredictable.)
\end{exercise}

\begin{exercise}[$\star\star$]\label{ex:chebsumapp}
Use Chebyshev's sum inequality to show that for $a,b,c>0$,
\[
\frac{a^3+b^3+c^3}3\ge\frac{a^2+b^2+c^2}3\cdot\frac{a+b+c}3\ge\Bigl(\frac{a+b+c}3\Bigr)^3,
\]
and more generally that $M_{k+1}^{k+1}\ge M_k^kM_1$ for the power means with equal weights and integer $k\ge1$.
\end{exercise}

\begin{exercise}[$\star\star\star$]\label{ex:schur}
\emph{Schur's inequality.} Show that for $a,b,c\ge0$,
\[
a(a-b)(a-c)+b(b-a)(b-c)+c(c-a)(c-b)\ge0,
\]
with equality iff $a=b=c$ or two of them are equal and the third is $0$. Deduce that $a^3+b^3+c^3+3abc\ge a^2(b+c)+b^2(c+a)+c^2(a+b)$.
\end{exercise}

\begin{challenge}
\begin{enumerate}[label=(\arabic*)]
\item \emph{The constant $e$ is best possible in Carleman's inequality.} Take $a_n=\frac1n$ for $n\le N$ and $a_n=0$ after, and compare both sides as $N\to\infty$. You will need Stirling's formula (Theorem~\ref{thm:stirlingC}) or the weaker asymptotic $\frac1n\ln(n!)=\ln n-1+o(1)$, obtained by integral comparison.
\item \emph{Hardy's inequality.} For non-negative $a_n$ with $\sum a_n^2<\infty$ and $A_n=a_1+\dots+a_n$, show that $\sum_{n\ge1}\bigl(\frac{A_n}n\bigr)^2\le4\sum_{n\ge1}a_n^2$.
\item \emph{The regular polygon wins.} Among all $n$-gons with $n\ge3$ distinct vertices on a given circle, listed in cyclic order, the regular one has the largest area.
\end{enumerate}
Notes at the end of the chapter.
\end{challenge}

\begin{exercise}[$\star\star$]\label{pb:stanley}
\emph{Counting ones.}
\begin{enumerate}[label=(\alph*)]
\item (Stanley.) Show that the total number of parts equal to $1$, over all partitions of $n$, equals the total number of
\emph{distinct} part sizes, over all partitions of $n$. Check it for $n=4$.
\item (Elder.) More generally, for $k\ge1$, show that the total number of parts equal to $k$ over all partitions of $n$ equals the
total number of part sizes that appear at least $k$ times, over all partitions of $n$.
\end{enumerate}
\end{exercise}

\section{Hints}

\begin{hint}{ex:8abc}
Bound each factor $a+b$ below by a product.
\end{hint}
\begin{hint}{ex:sqrtk}
You know $\sum(\sqrt k)^2$ exactly. Cauchy--Schwarz against the vector $(1,\dots,1)$.
\end{hint}
\begin{hint}{ex:intcs}
Write $1=\sqrt f\cdot\frac1{\sqrt f}$ and integrate.
\end{hint}
\begin{hint}{ex:euler}
AM--GM on $n+1$ well-chosen numbers: $n$ of them equal to $1+\frac1n$. What should the last one be so that the arithmetic mean is $1+\frac1{n+1}$?
\end{hint}
\begin{hint}{ex:nesbitt}
With $S=a+b+c$, $\frac a{b+c}=\frac S{S-a}-1$. Then use the Engel form.
\end{hint}
\begin{hint}{ex:triangle}
$A+B+C=\pi$ and $\sin$ is concave on $[0,\pi]$.
\end{hint}
\begin{hint}{ex:aabbcc}
Take logarithms, and use $(x-y)(\ln x-\ln y)\ge0$ for each pair.
\end{hint}
\begin{hint}{ex:holder}
(a) Weighted AM--GM with weights $\frac1p,\frac1q$. (b) Both sides are homogeneous: reduce to $\sum a_i^p=\sum b_i^q=1$.
\end{hint}
\begin{hint}{ex:carleman}
With $c_k=\frac{(k+1)^k}{k^{k-1}}$, compute $c_1c_2\cdots c_n$ and apply AM--GM to $c_1a_1,\dots,c_na_n$. Then swap the order of summation.
\end{hint}
\begin{hint}{ex:means3}
The middle inequality is AM--GM. For the left one, apply it to $\frac1a,\frac1b,\frac1c$. For the right one, square both sides and use $ab+bc+ca\le a^2+b^2+c^2$.
\end{hint}
\begin{hint}{ex:bernoulli}
(a) $x=\frac1n$, then $x=-\frac1{n^2}$. (b) $n=(1+x_n)^n\ge1+\binom n2x_n^2$ with $x_n\ge0$.
\end{hint}
\begin{hint}{ex:lagrangeid}
Expand $\sum_{i,j}a_i^2b_j^2-\sum_{i,j}a_ib_ia_jb_j$ and pair the terms $(i,j)$ and $(j,i)$. For $n=2$, write $a_1+ia_2$ and $b_1+ib_2$.
\end{hint}
\begin{hint}{ex:cauchyinduction}
(a) Split $2n$ numbers into two groups of $n$, apply $G_n\le A_n$ to each, then $G_2\le A_2$ to the two means. (b) With the extra number equal to the mean $m$ of the others, the arithmetic mean of all $n$ is still $m$. (c) Every $n$ is at most a power of $2$.
\end{hint}
\begin{hint}{ex:heron}
The three numbers $s-a,s-b,s-c$ are positive and have a fixed sum.
\end{hint}
\begin{hint}{ex:entropy}
Jensen for the concave function $\ln$ with weights $p_i$ at the points $\frac1{p_i}$.
\end{hint}
\begin{hint}{ex:chebsumapp}
Sort $a,b,c$; the lists $(a^2,b^2,c^2)$ and $(a,b,c)$ are sorted the same way. In general, apply the proposition to $(x_i^k)$ and $(x_i)$.
\end{hint}
\begin{hint}{ex:schur}
By symmetry assume $a\ge b\ge c$. Group the first two terms: $(a-b)\bigl[a(a-c)-b(b-c)\bigr]$, and check the sign of each piece.
\end{hint}

\begin{hint}{pb:stanley}
Count the partitions of $n$ that contain at least $j$ copies of the part $m$.
\end{hint}

\section{Solutions}

\begin{solution}{ex:8abc}
\idea The right side is a product, the left a product of sums: each sum $a+b$ should be bounded below by a product, and AM--GM for two numbers, $a+b\ge2\sqrt{ab}$, does exactly that. Check that the three bounds multiply to $8abc$.
\writeup By AM--GM, $a+b\ge2\sqrt{ab}$, $b+c\ge2\sqrt{bc}$, $c+a\ge2\sqrt{ca}$. All are positive, so we may multiply: the product is at least $8\sqrt{a^2b^2c^2}=8abc$.
\end{solution}

\begin{solution}{ex:sqrtk}
\idea We need an upper bound for a sum, in terms of something whose square we know: $\sum(\sqrt k)^2=\sum k$ is explicit. Cauchy--Schwarz against the constant vector $(1,\dots,1)$ bounds a sum by $\sqrt n$ times the root of the sum of squares.
\writeup Cauchy--Schwarz with $a_k=1$, $b_k=\sqrt k$: $\left(\sum\sqrt k\right)^2\le n\sum k=\frac{n^2(n+1)}2$. Take square roots.
\end{solution}

\begin{solution}{ex:intcs}
\idea A product of two integrals bounded below by a squared integral is Cauchy--Schwarz read backwards. Choose the two functions so that their product is $1$ and their squares are $f$ and $\frac1f$: they are $\sqrt f$ and $\frac1{\sqrt f}$.
\writeup Cauchy--Schwarz for integrals with $\sqrt f$ and $1/\sqrt f$: $1=\left(\int\sqrt f\cdot\frac1{\sqrt f}\right)^2\le\int f\int\frac1f$.
Equality holds if and only if $\sqrt f$ and $1/\sqrt f$ are proportional, $\sqrt f=\frac c{\sqrt f}$, that is, $f$ is constant.
\end{solution}

\begin{solution}{ex:euler}
\idea We want
\[
\left(1+\frac1n\right)^n<\left(1+\frac1{n+1}\right)^{n+1},
\]
that is, a geometric mean of $n+1$ numbers below their arithmetic mean $1+\frac1{n+1}$. A product $\left(1+\frac1n\right)^n$ of $n$ factors becomes a product of $n+1$ factors by adding a factor $1$. Check the arithmetic mean: it is the right one.
\writeup Apply AM--GM to the $n+1$ numbers $1,1+\frac1n,\dots,1+\frac1n$ ($n$ copies). Their mean is $\frac{1+(n+1)}{n+1}=1+\frac1{n+1}$, so
\[
\left(1+\frac1n\right)^{n/(n+1)}<1+\frac1{n+1},
\]
strictly since the numbers are not all equal. Raising to the power $n+1$ gives $\left(1+\frac1n\right)^n<\left(1+\frac1{n+1}\right)^{n+1}$.
\end{solution}

\begin{solution}{ex:nesbitt}
\idea The denominators $b+c$ differ from term to term, which is awkward. Writing $b+c=S-a$ with $S=a+b+c$ gives $\frac a{S-a}=\frac S{S-a}-1$, so we need a lower bound for $\sum\frac1{S-a}$: a sum of reciprocals with known total denominator, the textbook case of the Engel form.
\writeup Let $S=a+b+c$. Then $\frac a{b+c}=\frac S{S-a}-1$, so the left-hand side is $S\sum\frac1{S-a}-3$. By the Engel form of Cauchy--Schwarz
with $x_i=1$,
\[
\sum\frac1{S-a}\ge\frac9{(S-a)+(S-b)+(S-c)}=\frac9{2S}.
\]
Hence the sum is at least $\frac92-3=\frac32$, with equality if and only if $S-a=S-b=S-c$, i.e.\ $a=b=c$.
\end{solution}

\begin{solution}{ex:triangle}
\idea The constraint $A+B+C=\pi$ is linear and the quantity is a sum of values of one function: this is Jensen's setting, if $\sin$ has the right convexity on the range of angles. It is concave on $[0,\pi]$, which gives an upper bound at the mean angle $\frac\pi3$.
\writeup $\sin$ is concave on $[0,\pi]$ since $\sin''=-\sin\le0$ there (Proposition~\ref{prop:convex}). By Jensen with weights $\frac13$,
$\frac{\sin A+\sin B+\sin C}3\le\sin\frac{A+B+C}3=\sin\frac\pi3=\frac{\sqrt3}2$. Equality for the equilateral triangle.
\end{solution}

\begin{solution}{ex:aabbcc}
\idea Take logarithms: we want $3\sum a\ln a\ge(a+b+c)(\ln a+\ln b+\ln c)$. The left pairs each $x$ with its own $\ln x$, the right pairs all $x$ with all $\ln y$: ``matching'' beats ``mixing'', which is the rearrangement idea, and the hint is its two-term form.
\writeup $\ln$ is increasing, so $(x-y)(\ln x-\ln y)\ge0$, i.e.\ $x\ln x+y\ln y\ge x\ln y+y\ln x$. Sum over the three pairs $\{a,b\},\{b,c\},\{c,a\}$:
\[
2(a\ln a+b\ln b+c\ln c)\ge\sum_{x\ne y}x\ln y,
\]
where $(x,y)$ ranges over ordered pairs of distinct positions in the list $(a,b,c)$. Adding $a\ln a+b\ln b+c\ln c$ to both sides gives $3\sum a\ln a\ge\sum_{x,y}x\ln y=(a+b+c)(\ln a+\ln b+\ln c)$;
divide by $3$ and exponentiate. (This is also Chebyshev's sum inequality, Proposition~\ref{prop:chebsum}, for the lists $(a,b,c)$ and $(\ln a,\ln b,\ln c)$, which are sorted the same way.)
\end{solution}

\begin{solution}{ex:holder}
\idea (a) The right side is a weighted arithmetic mean of $a^p$ and $b^q$ with weights $\frac1p+\frac1q=1$, and the left side is the matching weighted geometric mean: weighted AM--GM. (b) H\"older is homogeneous (scaling $a$ scales both sides), so we may normalise to $\sum a_i^p=\sum b_i^q=1$, and then Young applied termwise sums to exactly $1$.
\writeup (a) If $a$ or $b$ is $0$ it is clear. Otherwise, by weighted AM--GM with weights $\frac1p,\frac1q$ applied to $a^p$ and $b^q$:
$ab=(a^p)^{1/p}(b^q)^{1/q}\le\frac{a^p}p+\frac{b^q}q$.

(b) Let $A=\left(\sum a_i^p\right)^{1/p}$ and $B=\left(\sum b_i^q\right)^{1/q}$; if $A$ or $B$ is $0$ all terms vanish. Apply (a) to $a_i/A$ and $b_i/B$ and sum:
\[
\frac{\sum a_ib_i}{AB}\le\frac1p\frac{\sum a_i^p}{A^p}+\frac1q\frac{\sum b_i^q}{B^q}=\frac1p+\frac1q=1.
\]
\end{solution}

\begin{solution}{ex:carleman}
\idea The geometric mean $(a_1\cdots a_n)^{1/n}$ is bounded by AM--GM, but the naive bound $\frac1n\sum_{k\le n}a_k$ gives, after summing over $n$, $\sum_ka_k\sum_{n\ge k}\frac1n=\infty$. The fix is to put weights $c_k$ inside before applying AM--GM and to divide by $(c_1\cdots c_n)^{1/n}$ outside: if this product is about $n^n$, the outer factor becomes $\frac1{n^2}$-like and the double sum converges. The choice $c_k=\frac{(k+1)^k}{k^{k-1}}$ makes the product telescope to $(n+1)^n$; then $\frac{c_k}k=\left(1+\frac1k\right)^k<e$ produces the constant.
\writeup The product telescopes:
\[\prod_{k=1}^nc_k=\frac{2^1\,3^2\cdots(n+1)^n}{1^0\,2^1\cdots n^{n-1}}=(n+1)^n.\]
So, by AM--GM,
\[
\begin{aligned}\left(a_1\cdots a_n\right)^{1/n}&=\frac{\left(c_1a_1\cdots c_na_n\right)^{1/n}}{n+1}\eqbr\le\frac1{n(n+1)}\sum_{k=1}^nc_ka_k.\end{aligned}
\]
The tangent to $\ln$ at $1$ lies above the graph, so $\ln(1+1/k)<1/k$ and $\bigl(1+\frac1k\bigr)^k<e$. Sum over $n\ge1$ and swap the sums by Proposition~\ref{prop:safe-sums}, which applies because every term is nonnegative. The inner sum telescopes, $\sum_{n\ge k}1/(n(n+1))=1/k$, and therefore
\[
\sum_{n=1}^\infty\left(a_1\cdots a_n\right)^{1/n}\le\sum_{k=1}^\infty\frac{c_k}ka_k=\sum_{k=1}^\infty\left(1+\frac1k\right)^ka_k\le e\sum_{k=1}^\infty a_k.
\]
For the value of the limit, compare integrals on $[1,1+1/k]$. There $1/t\ge k/(k+1)$, so
\[
\ln\left(1+\frac1k\right)=\int_1^{1+1/k}\frac1t\,\dd t\ge\int_1^{1+1/k}\frac{k}{k+1}\,\dd t=\frac{k}{k+1}\cdot\frac1k=\frac1{k+1}.
\]
Multiply by $k$: $k\ln(1+1/k)\ge k/(k+1)$. Exponentiate.
\[
e^{k/(k+1)}\le\left(1+\frac1k\right)^k<e.
\]
The lower end tends to $e$, so the middle term does too. Exercise~\ref{ex:euler} proves that the sequence is strictly increasing, which already gives existence of the limit once the upper bound is known. The integral supplies the value.
\end{solution}

\begin{solution}{ex:means3}
\idea Each step is AM--GM applied to a different list. For the harmonic mean, apply it to the reciprocals. For the quadratic mean, square both sides and use $ab+bc+ca\le a^2+b^2+c^2$ from the first section.
\writeup Middle: AM--GM for three numbers. Left: AM--GM for $\frac1a,\frac1b,\frac1c$ gives $\frac1{\sqrt[3]{abc}}\le\frac13\bigl(\frac1a+\frac1b+\frac1c\bigr)$; invert. Right: both sides are positive, so it is equivalent to $(a+b+c)^2\le3(a^2+b^2+c^2)$, i.e.\ $2(ab+bc+ca)\le2(a^2+b^2+c^2)$, which is the sum of $2ab\le a^2+b^2$ over the three pairs.
\end{solution}

\begin{solution}{ex:bernoulli}
\idea Bernoulli with the two obvious choices of $x$; then for $\sqrt[n]n$, the binomial theorem keeps a quadratic term that Bernoulli threw away, and that term is what forces $x_n\to0$.
\writeup (a) $x=\frac1n$: $(1+\frac1n)^n\ge1+n\cdot\frac1n=2$. $x=-\frac1{n^2}\ge-1$: $(1-\frac1{n^2})^n\ge1-\frac n{n^2}=1-\frac1n$.

(b) For $n\ge2$, $\sqrt[n]n>1$, so $x_n=\sqrt[n]n-1>0$ and, by the binomial theorem with all terms non-negative, $n=(1+x_n)^n\ge1+nx_n+\binom n2x_n^2\ge\frac{n(n-1)}2x_n^2$. Hence $x_n\le\sqrt{\frac2{n-1}}\to0$. (A derivative-free proof that $\ln n/n\to0$, in effect.)
\end{solution}

\begin{solution}{ex:lagrangeid}
\idea Both sides are sums over pairs of indices; expand the left side as a double sum, note that the diagonal $i=j$ cancels, and symmetrise the rest. The identity shows that the defect in Cauchy--Schwarz is a sum of squares, which is non-negative and vanishes exactly when all $2\times2$ determinants $a_ib_j-a_jb_i$ do, that is, when $a$ and $b$ are linearly dependent.
\writeup Expand the left side as
\[
\sum_{i,j}a_i^2b_j^2-\sum_{i,j}a_ib_ia_jb_j
=\sum_{i,j}\bigl(a_i^2b_j^2-a_ib_ja_jb_i\bigr).
\] The terms with $i=j$ vanish. Pair $(i,j)$ with $(j,i)$ for $i<j$: their contribution is $a_i^2b_j^2+a_j^2b_i^2-2a_ib_ja_jb_i=(a_ib_j-a_jb_i)^2$. So the left side is $\sum_{i<j}(a_ib_j-a_jb_i)^2\ge0$, which is Cauchy--Schwarz; it is $0$ iff $a_ib_j=a_jb_i$ for all $i,j$, iff $a$ and $b$ are linearly dependent (if $b\ne0$, say $b_j\ne0$, then $a_i=\frac{a_j}{b_j}b_i$ for all $i$).

For $n=2$: $(a_1^2+a_2^2)(b_1^2+b_2^2)=(a_1b_1+a_2b_2)^2+(a_1b_2-a_2b_1)^2$, which is $|z|^2|w|^2=|z\bar w|^2$ for $z=a_1+ia_2$, $w=b_1+ib_2$: the identity of Brahmagupta from Chapter~\ref{ch:complex}, and the reason a product of two sums of two squares is a sum of two squares.
\end{solution}

\begin{solution}{ex:cauchyinduction}
\idea Doubling is natural: $2n$ numbers split into two halves, each half has a geometric mean below its arithmetic mean, and the two-number case combines them. Going down is the clever step: adding the mean of the others as an extra number does not change the arithmetic mean, so $G_n\le A_n$ for this special list compares $G_{n-1}$ with $A_{n-1}$. Together, the two steps reach every $n$ from the powers of $2$.
\writeup (a) $G_2\le A_2$ is $\sqrt{xy}\le\frac{x+y}2$, from $(\sqrt x-\sqrt y)^2\ge0$. Given $x_1,\dots,x_{2n}$, let $g,g'$ and $m,m'$ be the geometric and arithmetic means of the first $n$ and of the last $n$. Then $G_{2n}=\sqrt{gg'}\le\frac{g+g'}2\le\frac{m+m'}2=A_{2n}$, using $G_2\le A_2$ and then $g\le m$, $g'\le m'$.

(b) Let $m=\frac{x_1+\dots+x_{n-1}}{n-1}$ and apply $G_n\le A_n$ to $x_1,\dots,x_{n-1},m$: the arithmetic mean of these $n$ numbers is $\frac{(n-1)m+m}n=m$, so $(x_1\cdots x_{n-1}m)^{1/n}\le m$, i.e.\ $x_1\cdots x_{n-1}\le m^{n-1}$, which is $G_{n-1}\le A_{n-1}$.

(c) By (a) and induction, $G_{2^k}\le A_{2^k}$ for all $k$. Given $n$, choose $k$ with $2^k\ge n$ and go down from $2^k$ to $n$ by (b).
\end{solution}

\begin{solution}{ex:heron}
\idea The area is a product of three factors with a fixed sum, $(s-a)+(s-b)+(s-c)=s$: AM--GM bounds the product by the cube of the mean, with equality when the three factors are equal, that is, $a=b=c$. (The factors are positive by the triangle inequality: $s-a=\frac{b+c-a}2>0$.)
\writeup By AM--GM, $(s-a)(s-b)(s-c)\le\Bigl(\frac{(s-a)+(s-b)+(s-c)}3\Bigr)^3=\frac{s^3}{27}$, with equality iff $s-a=s-b=s-c$, i.e.\ $a=b=c=\frac{2s}3$. So the area is at most $\sqrt{s\cdot\frac{s^3}{27}}=\frac{s^2}{3\sqrt3}$, the area of the equilateral triangle of side $\frac{2s}3$ (check: $\frac{\sqrt3}4\bigl(\frac{2s}3\bigr)^2=\frac{s^2}{3\sqrt3}$).
\end{solution}

\begin{solution}{ex:entropy}
\idea $H$ is a weighted average of the values $\ln\frac1{p_i}$ with weights $p_i$: Jensen for the concave $\ln$ moves the average inside, where it becomes $\ln\sum p_i\cdot\frac1{p_i}=\ln n$.
\writeup Jensen for the concave function $\ln$ (Theorem~\ref{thm:jensen} applied to $-\ln$) with weights $p_i$ at the points $\frac1{p_i}$:
\[
H=\sum_ip_i\ln\frac1{p_i}\le\ln\Bigl(\sum_ip_i\cdot\frac1{p_i}\Bigr)=\ln n.
\]
Since $\ln''<0$, equality forces all the points $\frac1{p_i}$ equal, so $p_i=\frac1n$.
\end{solution}

\begin{solution}{ex:chebsumapp}
\idea Chebyshev's inequality compares the mean of products with the product of means, for similarly sorted lists. Powers of the same positive numbers are always sorted the same way, so $(x_i^k)$ and $(x_i)$ qualify, and the product $x_i^k\cdot x_i=x_i^{k+1}$ is the next power.
\writeup Sort so that $a\le b\le c$; then $a^2\le b^2\le c^2$ and Proposition~\ref{prop:chebsum} gives $\frac{a^3+b^3+c^3}3\ge\frac{a^2+b^2+c^2}3\cdot\frac{a+b+c}3$. Apply it again to $(a,b,c)$ and $(a,b,c)$: $\frac{a^2+b^2+c^2}3\ge\bigl(\frac{a+b+c}3\bigr)^2$, which gives the second inequality. In general, with $x_1\le\dots\le x_n$, the lists $(x_i^k)$ and $(x_i)$ are sorted the same way, so $\frac1n\sum x_i^{k+1}\ge\frac1n\sum x_i^k\cdot\frac1n\sum x_i$, which is $M_{k+1}^{k+1}\ge M_k^kM_1$. (Iterating, $M_k^k\ge M_1^k$: the power mean inequality $M_1\le M_k$ for integers $k$, without Jensen.)
\end{solution}

\begin{solution}{ex:schur}
\idea The expression is symmetric, so we may order the variables, and then the trick is to pair the two largest terms: their sum factors through $a-b\ge0$ and a bracket that is $\ge0$ because $a\ge b$ and $a-c\ge b-c\ge0$; the third term is non-negative on its own. The second inequality is just the first one expanded.
\writeup By symmetry assume $a\ge b\ge c\ge0$. Then
\[
a(a-b)(a-c)+b(b-a)(b-c)=(a-b)\bigl[a(a-c)-b(b-c)\bigr]\ge0,
\]
because $a-b\ge0$, $a\ge b\ge0$ and $a-c\ge b-c\ge0$ give $a(a-c)\ge b(b-c)$. And $c(c-a)(c-b)=c(a-c)(b-c)\ge0$. Adding proves the inequality. Equality needs both pieces to vanish. If $a>b$, the bracket is $>0$: indeed $a(a-c)>b(a-c)\ge b(b-c)$ since $a-c>0$. So the first piece vanishes iff $a=b$. The second vanishes iff $c=0$ or $c=b$ (then $a=b=c$) or $c=a$ (then $a=b=c$ again). Combining: $a=b=c$, or $a=b$ and $c=0$; by symmetry, equality holds iff all three are equal, or two are equal and the third is $0$.

Expanding, $\sum a(a-b)(a-c)=\sum a^3-\sum a^2(b+c)+3abc$ (each term $a(a-b)(a-c)=a^3-a^2b-a^2c+abc$), so the inequality reads $a^3+b^3+c^3+3abc\ge a^2(b+c)+b^2(c+a)+c^2(a+b)$.
\end{solution}

\subsection*{Notes on the challenges}

(1) With $a_n=\frac1n$ for $n\le N$: the right side is $e\sum_{n\le N}\frac1n=e\ln N+O(1)$ (Chapter~\ref{ch:sequences}). On the left, $(a_1\cdots a_n)^{1/n}=(n!)^{-1/n}$, and Stirling's formula gives $(n!)^{1/n}\sim\frac ne$, so the left side is $\sum_{n\le N}\frac en(1+o(1))=e\ln N+o(\ln N)$. The ratio of the two sides tends to $1$: no constant smaller than $e$ works. Numerical evaluation shows the ratio of the left side to $\sum_{n\le N}a_n$ creeping up: $1.52$ for $N=10$, $1.94$ for $N=100$, $2.29$ for $N=10^4$; the convergence is logarithmically slow, which is why the numerical evidence alone is unconvincing.

(2) Hardy's inequality (1920, constant $4=\bigl(\frac p{p-1}\bigr)^p$ for $p=2$) can be proved with Cauchy--Schwarz and a summation by parts; Steele's book, Chapter~11, has a guided proof, and the constant is again best possible, but the test sequence is different: take $a_n=n^{-1/2}$ for $1\le n\le N$ and $a_n=0$ afterwards. For $n\le N$, integral comparison gives $A_n=2\sqrt n+O(1)$, so
\[
\sum_{n=1}^N\left(\frac{A_n}{n}\right)^2=4\ln N+O(1),
\qquad
\sum_{n=1}^N a_n^2=\ln N+O(1).
\]
The terms with $n>N$ on the left are non-negative, so no constant below $4$ can work for all $N$.

(3) An $n$-gon inscribed in a circle of radius $R$ whose centre lies inside it has area $\frac{R^2}2\sum_{i=1}^n\sin\theta_i$, where $\theta_i\in\,]0,\pi[$ are the central angles, with $\sum\theta_i=2\pi$. Jensen for the concave $\sin$ gives $\sum\sin\theta_i\le n\sin\frac{2\pi}n$, with equality for the regular polygon. If the centre is outside (or on a side), the same formula holds with signed angles: one central angle $\theta_1$ is now in $[\pi,2\pi[$ and contributes $\sin\theta_1\le0$, while the others sum to $2\pi-\theta_1\le\pi$, so Jensen gives $\sum\sin\theta_i\le(n-1)\sin\frac{\pi}{n-1}<n\sin\frac{2\pi}n$ (for $n=3$, compare $2$ with $3\sqrt3/2$; for $n\ge4$, use $(n-1)\sin\frac\pi{n-1}\le\pi<4\le n\sin\frac{2\pi}n$, the last inequality following from the decrease of $\sin x/x$ on $]0,\pi/2]$): the area is smaller than in the first case.

\begin{solution}{pb:stanley}
\idea Count the same pairs in two ways: a partition together with ``a $j$-th copy of the part $m$''.

\writeup (b) contains (a) as the case $k=1$, so we prove (b). Fix $m\ge1$ and $j\ge1$. Removing $j$ copies of $m$ is a bijection
from the partitions of $n$ with at least $j$ copies of $m$ onto the partitions of $n-jm$, so there are $p(n-jm)$ of them (with
$p$ of a negative number equal to $0$).

The total number of parts equal to $k$ over all partitions of $n$ is $\sum_\lambda m_k(\lambda)$, where $m_k$ is the multiplicity of $k$.
Since $m_k(\lambda)=\#\{j\ge1:m_k(\lambda)\ge j\}$, this total is $\sum_{j\ge1}\#\{\lambda:m_k(\lambda)\ge j\}=\sum_{j\ge1}p(n-jk)$.

The total number of part sizes appearing at least $k$ times is $\sum_{m\ge1}\#\{\lambda:m_m(\lambda)\ge k\}=\sum_{m\ge1}p(n-km)$. The two
sums are the same sum, with the summation variable renamed.

For $n=4$ the partitions are $4$, $3+1$, $2+2$, $2+1+1$, $1+1+1+1$. They contain $0+1+0+2+4=7$ ones and $1+2+1+2+1=7$
distinct part sizes.

\tools Chapter~\ref{ch:counting} (counting one set of pairs in two ways).
\leads (a) was found by Richard Stanley, and (b) by Paul Elder in 1984; Honsberger's \emph{Mathematical Gems~III} tells the
story. The common value in (a) is $\sum_{j\ge1}p(n-j)$, so the average number of $1$'s in a partition of $n$ is
$\frac{p(0)+\dots+p(n-1)}{p(n)}$, which grows like $\frac{\sqrt{6n}}\pi$ by the Hardy--Ramanujan asymptotics.
\end{solution}

\section*{Sources and further reading}

\begin{itemize}\raggedright
\item J.~M.~Steele, \emph{The Cauchy--Schwarz Master Class}, Cambridge University Press, 2004. Inequalities taught by proofs and problems. Chapters~1 and~2 treat Cauchy--Schwarz and AM--GM, and the historical notes are the source for the attributions to Cauchy in $1821$, Bunyakovsky in $1859$ and Schwarz in $1885$.
\item G.~H.~Hardy, J.~E.~Littlewood and G.~P\'olya, \emph{Inequalities}, 2nd ed., Cambridge University Press, 1952. The classic reference: convexity, AM--GM, H\"older, power means, Chebyshev's sum inequality, rearrangement (Chapter~10) and Carleman.
\item J.~L.~W.~V.~Jensen, ``Sur les fonctions convexes et les in\'egalit\'es entre les valeurs moyennes'', \emph{Acta Mathematica} 30 (1906), 175--193. The paper of the telephone engineer; the biographical facts above are from the MacTutor biography.
\item A.-L.~Cauchy, \emph{Cours d'analyse de l'\'Ecole royale polytechnique}, Paris, 1821, Note~II. The comparison of arithmetic and geometric means, rebuilt in Exercise~\ref{ex:cauchyinduction}, and the sum form of Theorem~\ref{thm:cs}.
\item V.~Bouniakowsky, ``Sur quelques in\'egalit\'es concernant les int\'egrales ordinaires et les int\'egrales aux diff\'erences finies'', \emph{M\'emoires de l'Acad\'emie imp\'eriale des sciences de St.-P\'etersbourg}, VIIe s\'erie, tome~1, no.~9, 1859.
\item H.~A.~Schwarz, ``\"Uber ein die Fl\"achen kleinsten Fl\"acheninhalts betreffendes Problem der Variationsrechnung'', \emph{Acta Societatis Scientiarum Fennicae} 15 (1885), 315--362. The quadratic-polynomial proof of the integral form.
\item J.~Bernoulli, \emph{Positiones arithmeticae de seriebus infinitis}, Basel, 1689. The inequality of Proposition~\ref{prop:bernoulli}, under the name it still carries.
\item I.~Barrow, \emph{Lectiones geometricae}, London, 1670, Lecture~7. The same comparison, earlier than Bernoulli's tract.
\item G.~P\'olya, \emph{Proof of an inequality}, Proc.\ London Math.\ Soc.\ (2) 24 (1926), p.~lvii. The weighted AM--GM proof of Carleman's inequality used in Exercise~\ref{ex:carleman}.
\item J.~Duncan and C.~M.~McGregor, \emph{Carleman's inequality}, Amer.\ Math.\ Monthly 110 (2003), 424--431. A survey of proofs and of why $e$ is optimal.
\item A.~Engel, \emph{Problem-Solving Strategies}, Springer, 1998, Chapter~7. Many olympiad inequalities in the style of Exercises~\ref{ex:8abc}--\ref{ex:aabbcc} and~\ref{ex:schur}.
\item T.~Carleman, \emph{Sur les fonctions quasi-analytiques}, Conf\'erences faites au cinqui\`eme congr\`es des math\'ematiciens scandinaves, Helsinki, 1923, 181--196.
\item H.~Minkowski, \emph{Geometrie der Zahlen}, Teubner, Leipzig, 1896. Where the $p$-norms and their triangle inequality were born, in the service of number theory.
\end{itemize}
